% !TeX program = xelatex
% ============================================================
% pdp_assignment/pdp_assignment.tex — «Индивидуальное задание на
% преддипломную практику» (ИЗ) по официальной форме
% «0-1-Форма индивидуального задания ПДП».
%
% Блок «Заявление на ПДП» в СмартЛМС (дедлайн ~12 апреля): ИЗ и Рабочий
% план с подписями студента и руководителя(-ей) практики. Порядок
% подписания: студент подписывает САМ → (при двух руководителях) сначала
% руководитель от профильной организации (соруководитель) → затем
% формальный руководитель от НИУ ВШЭ. Официальная форма русская — документ
% один для всех языков проекта. Данные — из настроек (\hse…), жёлтые
% \hseFill заполняются вручную.
% ============================================================
\documentclass[12pt,a4paper]{extarticle}
\input{../common/preamble}
\newcommand{\docname}{Индивидуальное задание}
\newcommand{\doccode}{}
\newcommand{\doccodelu}{}
\input{../common/commands}
\begin{document}
\pagestyle{empty}
\sloppy
\setlength{\parindent}{0pt}
\setlength{\parskip}{5pt}

\begin{center}
Федеральное государственное автономное образовательное учреждение\\
высшего образования\\
«Национальный исследовательский университет «Высшая школа экономики»

\vspace{0.9em}
\textbf{ИНДИВИДУАЛЬНОЕ ЗАДАНИЕ, ВЫПОЛНЯЕМОЕ В ПЕРИОД ПРАКТИКИ}

студентом \hseCourse{} курса очной формы обучения
\hseFill{ФИО в творительном падеже (Ивановым Иваном Ивановичем)}
\end{center}

\vspace{0.6em}
\begin{tabularx}{\textwidth}{@{}p{5.6cm}X@{}}
образовательной программы & «\hseSpecialization» \\
уровень & \hseDegreeLevel \\
по направлению & \hseProgramCode\ «\hseSpecialization» \\
факультета & \hseFaculty \\
Вид практики & производственная \\
Тип практики & преддипломная \\
Срок прохождения практики & с \hseFill{01 апреля \hseYear{} г.} по \hseFill{30 апреля \hseYear{} г.} \\
\end{tabularx}

\vspace{0.9em}
\textbf{Цель прохождения практики}

Преддипломная практика проводится с целью завершения подготовки выпускной
квалификационной работы (ВКР) по утверждённой теме.

\textbf{Задачи практики}
((* if project.kind == 'research' *))
\begin{enumerate}[leftmargin=1.4em, itemsep=2pt]
  \item Завершение исследования: проведение вычислительных экспериментов,
        получение заявленных результатов и проверка гипотез.
  \item Подготовка завершённого текста ВКР со всеми необходимыми разделами
        и приложениями.
  \item Подготовка презентации для отчёта о результатах преддипломной
        практики, предзащиты выпускной квалификационной работы.
\end{enumerate}
((* else *))
\begin{enumerate}[leftmargin=1.4em, itemsep=2pt]
  \item Завершение разработки и испытаний программного проекта в соответствии
        с утверждённым техническим заданием.
  \item Подготовка программной документации на разработанный программный проект.
  \item Подготовка завершённого текста ВКР со всеми необходимыми разделами
        и приложениями.
  \item Подготовка презентации для отчёта о результатах преддипломной
        практики, предзащиты выпускной квалификационной работы.
\end{enumerate}
((* endif *))

\textbf{Содержание практики}
\begin{enumerate}[leftmargin=1.4em, itemsep=2pt]
  \item \hseFill{работа 1 (например: доработка текста ВКР по замечаниям руководителя)};
  \item \hseFill{работа 2 (например: утверждение программной документации)};
  \item \hseFill{работа 3 (например: подготовка презентации и иллюстраций к защите)}.
\end{enumerate}

\textbf{Планируемые результаты}
\begin{enumerate}[leftmargin=1.4em, itemsep=2pt]
  \item \hseFill{финальная версия текста ВКР};
  \item \hseFill{готовая презентация к защите};
  \item \hseFill{работоспособный продукт / результаты экспериментов};
  \item \hseFill{полный пакет программной документации}.
\end{enumerate}

\vspace{1.2em}
Руководитель практики от НИУ ВШЭ:

\vspace{0.9em}
\begin{tabularx}{\textwidth}{@{}Xp{3.2cm}p{4.2cm}@{}}
\hseSupervisorTitle & \hrulefill & \hseSupervisorName \\
{\scriptsize (должность)} & {\scriptsize (подпись)} & {\scriptsize (фамилия, инициалы)} \\
\end{tabularx}

\ifhseHasCoSupervisor
\vspace{0.8em}
\textbf{СОГЛАСОВАНО}

Руководитель практики от профильной организации:

\vspace{0.9em}
\begin{tabularx}{\textwidth}{@{}Xp{3.2cm}p{4.2cm}@{}}
\hseFill{должность в организации} & \hrulefill & \hseCoSupervisorName \\
{\scriptsize (должность)} & {\scriptsize (подпись)} & {\scriptsize (фамилия, инициалы)} \\
\end{tabularx}
\fi

\vspace{1.2em}
Задание принято к исполнению \hseFill{«\_\_» апреля \hseYear{} г.}

\vspace{0.9em}
\begin{tabularx}{\textwidth}{@{}p{4cm}p{3.2cm}X@{}}
Студент & \hrulefill & \hseAuthorName \\
 & {\scriptsize (подпись)} & {\scriptsize (фамилия, инициалы)} \\
\end{tabularx}

\end{document}
